\chapter{1977年江苏徐州市、徐州地区卷}

\begin{problem}
计算：\( (-16)\div\left( -2\frac{2}{3} \right)^2+\left( \frac{1}{3} \right)^{-2}\times\left( \frac{3}{8} \right)^0-\left( \frac{9}{16} \right)^{\frac{1}{2}} \) 。

\end{problem}

\begin{problem}
（1）把 \( \frac{1}{\sqrt{3}-2} \) 的分母有理化，（2）已知 \( x=\sqrt{2}+1 \)，求 \( x^2-2x+1 \) 的值。

\end{problem}

\begin{problem}
求下列三角函数值：（1）\( \cos 150^\circ \)；（2）\( \operatorname{tg}\left( -\frac{\pi}{6} \right) \) 。

\end{problem}

\begin{problem}
解方程组：\( \left\{ \begin{aligned} 7x+6y&=-4 \\ 2x-3y&=13 \end{aligned} \right. \) 。

\end{problem}

\begin{problem}
计算：\( \frac{1}{x^2-1}+\frac{1}{x^2+4x+3} \) 。

\end{problem}

\begin{problem}
一船在 \( A \) 点测得灯塔 \( C \) 在其北偏东 \( 45^\circ \) 的方向，此船以每小时18海浬的速度向南航行半小时至 \( B \) 后，又测得灯塔在其北偏东 \( 15^\circ \) 的方向，求此时船与灯塔的距离 \( BC \) 。

\bitmapfigure[max width=2.8cm,max totalheight=5.0cm]{img/1977/jiangsu_xuzhou_city_region/q6_fig1.png}

\end{problem}

\begin{problem}
某生产队准备在一边靠墙的地方建一个长方形的猪圈，现有一堆旧砖，可以垒成20米长的围墙（门在内）。问应怎样选取猪圈的宽才能使垒成的猪圈面积最大？最大面积是多少？

\bitmapfigure[max width=6.0cm,max totalheight=2.4cm]{img/1977/jiangsu_xuzhou_city_region/q7_fig1.png}

\end{problem}

\begin{problem}
已知：\( \lg 2=0.3010 \)，\( \lg 3=0.4771 \)。计算从1到10之间（除7以外）所有整数的对数值。

\end{problem}

\begin{problem}
已知顶点在原点的抛物线，关于 \( y \) 轴对称，并且经过 \( P(6,-2) \) 点。（1）求抛物线的方程；（2）画出它的大致图形；（3）写出它的焦点坐标和准线方程。

\end{problem}

\begin{problem}
已知：\( a \)，15，\( b \) 三个数成等差数列，\( a \)，9，\( b \) 三个数成等比数列。求 \( a \)、\( b \) 两个数。
\end{problem}

